(1) 取平面法向为 \(\hat y\)，由高斯柱面与对称性，静止系有
\[
\boldsymbol E'_\pm=\pm\frac{\sigma_0}{2\epsilon_0}\hat y,\qquad \boldsymbol B'=0,
\]
正负号对应 y>0 与 y<0。平面的无限大假设使电场与距平面的远近无关。

(2) 令 S 相对于 S′ 沿 +x 以速度 v 运动，因此在 S 中电荷平面沿 −x 方向移动。按从 S′ 到 S 的场变换，垂直速度的场为 \(E_\perp=\gamma E'_\perp\)，且 \(\boldsymbol B_\perp=-\gamma\,\boldsymbol v\times\boldsymbol E'/c^2\)（这里 \(\boldsymbol v=+v\hat x\) 指坐标系 S 相对 S′ 的速度）。所以
\[
\boldsymbol E_\pm=\pm\frac{\gamma\sigma_0}{2\epsilon_0}\hat y,\quad
\boldsymbol B_\pm=\mp\frac{\gamma v\sigma_0}{2\epsilon_0c^2}\hat z.
\]
由于电荷面内长度发生 Lorentz 收缩，S 中面密度为 \(\sigma=\gamma\sigma_0\)，与法向电场跳跃一致。

(3) S 中电荷流动方向为 −x，因此表面电流 \(\boldsymbol K=-\sigma v\hat x\)。无限电流平面的边界磁场为 \(\boldsymbol B_+=\mu_0(\boldsymbol K\times\hat y)/2=-\mu_0\sigma v\hat z/2\)，因为 \(\hat x\times\hat y=\hat z\)。借助 \(\mu_0\epsilon_0c^2=1\)，正好与上一部分 B+ 的结果一致。

(4) E 沿 y、B 沿 z，因此 \(\boldsymbol E\cdot\boldsymbol B=0\)。此外
\[
E^2-c^2B^2=
\frac{\gamma^2\sigma_0^2}{4\epsilon_0^2}(1-v^2/c^2)
=\frac{\sigma_0^2}{4\epsilon_0^2}.
\]
这与静止系 \(E'^2-c^2B'^2=E'^2\) 一致，展示两个洛伦兹标量不变量与实际表面电流场方程的一致性。
